\chapter*{贯穿案例：从结果 21 开始}
\addcontentsline{toc}{chapter}{贯穿案例：从结果 21 开始}

先把本册要追踪的第一次运行写成一句可以判定真假的要求：

\begin{quote}
机器从规定的复位状态出发，依次读取 \(7,-2,11,5\)，在地址
\code{0x00102020} 写入 32 位整数 \(21\)，然后回到启动程序。
\end{quote}

这句话故意没有使用“进程”“系统调用”或“文件”。当前机器还没有这些对象，只有处理器、
启动存储、RAM、地址互连、时钟、复位电路，以及在复位保持期间帮助写入 RAM 的外部夹具。
操作人员必须先把一段 32 字节的任务代码、四个输入整数、结果哨兵和返回位置放到确切地址，
再释放复位。处理器能够做的只是取指、发起存储事务、更新寄存器和按照指令改变控制位置。

为避免“差不多在这附近”掩盖错误，第一次运行采用固定布局：

\begin{longtable}{|L{0.31\linewidth}|L{0.22\linewidth}|L{0.35\linewidth}|}
\hline
\rowcolor{BookTableHead}
地址或范围 & 运行前的内容 & 成功运行后的可见事实 \\ \hline
\code{0x00100000--0x0010001f} & 八条定长求和指令 & 字节保持不变 \\ \hline
\code{0x00102000--0x0010200f} & \(7,-2,11,5\) 的小端表示 & 四个输入保持不变 \\ \hline
\code{0x00102020--0x00102023} & 哨兵 \code{ff ff ff ff} & 变为 \code{15 00 00 00} \\ \hline
\code{0x001efffc--0x001effff} & 启动程序返回地址 \code{0x00000300} & 被 \code{RET} 读取，SP 前移 4 \\ \hline
\end{longtable}

地址表只说明字节放在哪里，并没有产生保护。RAM 不知道某个范围叫“代码”，也不知道
\code{0x00102020} 是“结果”；处理器沿取指路径读取代码，沿装入路径解释输入，沿返回路径
把栈顶四字节写入 PC，这些字节才在各自的使用时刻获得含义。如果 PC 被破坏后指向输入区，
机器会尝试把整数 \(7\) 的字节当作指令；如果存储指令写向代码区，RAM 也不会因为“那里
本来是程序”而自动拒绝。

本次运行有三个不同的完成边界。RAM 接受最后一笔结果写入时，只能说明数值 \(21\) 已经
对处理器可见；任务执行 \code{RET} 并成功取得返回位置时，才能说明控制权已经离开任务；
启动程序确认返回状态且不再修改结果时，外部观察者才拥有一份稳定结果。把三者合成一句
“运行完成”，会使结果写出以后失控、返回地址损坏和观察过早变得无法区分。

后续章节不会更换这套提问方式。机器能够自动接收新任务以后，我们仍追问字节由谁拥有、
何时可执行；出现第二项任务以后，我们仍追问哪一份寄存器和栈属于谁；时钟能够打断任务
以后，我们仍追问保存到哪里才算可以恢复。对象会增加，地址和寄存器会换成 x86-64 的
真实形态，但每一次演进都要回到相同的六个问题：输入与前提是什么，对象由谁拥有，数据与
控制怎样移动，哪一步提交，外部看见什么，失败后留下什么。

\begin{problemframe}
\textbf{第一章的唯一目标。} 先证明这一项任务可以从复位开始得到 \(21\) 并合作返回。
机器怎样在运行中接收另一份程序、怎样保存第二份现场、怎样从不合作的任务手中取回
处理器，分别留给后续问题；第一章不借用尚未建立的机制。
\end{problemframe}
